\documentclass[../../main.tex]{subfiles}

\begin{document}

\section{可数性与可分性}\label{section:ch11-可数性可分性}

在度量空间中我们已定义了序列收敛的含义. 以下是序列的收敛性在拓扑空间的自然推广.

\begin{definition}[序列收敛]\label{definition:ch11-序列收敛}
拓扑空间 $X$ 中的序列 $\{x_{n}\}$ 称为收敛到点 $x\in X$, 若对 $x$ 的每个邻域 $U$, 存在指标 $N$ 使得若 $n\geqslant N$, 则 $x_{n}$ 属于 $U$. 点 $x$ 称为该序列的一个极限.
\end{definition}

在度量空间中, 一个序列不可能收敛于两个不同的点. 在一般的拓扑空间中, 一个序列可能收敛于两个不同的点. 例如, 对于集合上的平凡拓扑, 每个序列收敛到每个点. 对于 Hausdorff 空间, 序列有唯一的极限.

\begin{definition}[第一与第二可数]\label{definition:ch11-第一第二可数}
拓扑空间 $X$ 称为第一可数的, 若在每一点存在可数基. 空间 $X$ 称为第二可数的, 若存在该拓扑的可数基.
\end{definition}

显然第二可数空间是第一可数的.

\begin{example}\label{example:ch11-度量第一可数}
每个度量空间 $X$ 是第一可数的. 开球的可数族 $\{B(x,1/n)\}_{n=1}^{\infty}$ 是由度量诱导的拓扑在 $x$ 的基.
\end{example}

我们将以下命题的证明留作练习.

\begin{proposition}\label{proposition:ch11-第一可数闭包}
令 $X$ 为第一可数拓扑空间. 对于 $X$ 的子集 $E$, 点 $x\in X$ 是 $E$ 的闭包点当且仅当它是 $E$ 中序列的极限. 因此 $X$ 的子集 $E$ 是闭的当且仅当只要 $E$ 中的点收敛到 $x\in X$, 点 $x$ 属于 $E$.
\end{proposition}

在不是第一可数的拓扑空间, 可能一个点是闭包点但没有序列收敛到该点(见习题 22).

\begin{definition}[稠密与可分]\label{definition:ch11-稠密可分}
拓扑空间 $X$ 的子集 $E$ 称为在 $X$ 中稠密, 若 $X$ 的每个开集包含 $E$ 的点. 我们称 $X$ 为可分的, 若它有可数的稠密子集.
\end{definition}

显然 $E$ 在 $X$ 中稠密当且仅当 $X$ 中的每个点是 $E$ 的闭包点, 即 $\overline{E}=X$.

在第 9 章, 我们证明了一个度量空间是第二可数的当且仅当它是可分的. 在一般的拓扑空间, 第二可数空间是可分的, 但一个可分空间, 即使是第一可数的, 也可以不是第二可数的(见习题 21).

拓扑空间称为可度量化的, 若该拓扑可被度量诱导. 不是每个拓扑都可被度量诱导. 事实上, 我们已看到度量空间是正规的, 因此当然地, 多于一点的集合上的平凡拓扑是不可度量化的. 自然要问是否可能识别那些可度量化的拓扑空间. 谈到这一点我们意味着用拓扑的开集叙述拓扑可由度量诱导的充分与必要准则. 存在这样的准则. 在拓扑空间 $X$ 是第二可数的情形, 存在可度量化的简单的充分与必要准则.

\begin{theorem}[Urysohn 度量化定理]\label{theorem:ch11-urysohn度量化}
令 $X$ 为第二可数拓扑空间. 则 $X$ 是可度量化的当且仅当它是正规的.
\end{theorem}

我们已经证明一个度量空间是正规的. 对第二可数拓扑空间, 逆命题也成立. 我们把它的证明推迟到下一章.

\begin{exercise}\label{exercise:ch11-17}
拓扑空间称为 Lindel\"of 空间或有 Lindel\"of 性质, 若 $X$ 的每个开覆盖有可数子覆盖. 证明: 若 $X$ 是第二可数的, 则它是 Lindel\"of 的.
\end{exercise}

\begin{exercise}\label{exercise:ch11-18}
令 $X$ 为不可数点集, 且令 $\mathcal{T}$ 为 $\varnothing$ 与 $X$ 的所有具有有限补的子集组成的集族. 证明 $\mathcal{T}$ 是 $X$ 的拓扑且空间 $(X,\mathcal{T})$ 不是第一可数的.
\end{exercise}

\begin{exercise}\label{exercise:ch11-19}
证明第二可数空间是可分的且第二可数空间的每个子空间是第二可数的.
\end{exercise}

\begin{exercise}\label{exercise:ch11-20}
证明 Moore 平面是可分的(见习题 10). 证明 Moore 平面的子空间 $\mathbf{R}\times\{0\}$ 是不可分的. 得出 Moore 平面不是可度量化的且不是第二可数的.
\end{exercise}

\begin{exercise}\label{exercise:ch11-21}
证明 Sorgenfrey 直线是第一可数的但不是第二可数的. 得出 Sorgenfrey 直线不是可度量化的.
\end{exercise}

\begin{exercise}\label{exercise:ch11-22}
令 $X_{1}=\mathbf{N}\times\mathbf{N}$, 其中 $\mathbf{N}$ 表示自然数集且取 $X=X_{1}\cup\{\omega\}$, 其中 $\omega$ 不属于 $X_{1}$. 对每个自然数列 $s=\{m_{k}\}$ 与自然数 $n$, 定义
\begin{align*}
B_{s,n}=\{\omega\}\cup\{(j,k):\text{对所有 }k\geqslant n,j\geqslant m_{k}\}.
\end{align*}
\begin{enumerate}[(i)]
\item 证明集合 $B_{s,n}$ 与单点集 $\{(j,k)\}$ 一起构成 $X$ 上的拓扑的基.
\item 证明 $\omega$ 是 $X_{1}$ 的闭包点, 即使不存在 $X_{1}$ 中的序列 $\{x_{n}\}$ 收敛到 $\omega$.
\item 证明空间 $X$ 是可分的但不是第一可数的, 从而不是第二可数的.
\item $X$ 是 Lindel\"of 空间吗?
\end{enumerate}
\end{exercise}

\end{document}